\documentclass[11pt]{article}

\usepackage[utf8]{inputenc}
\usepackage[T1]{fontenc}
\usepackage{amsmath,amssymb,amsthm}
\usepackage{geometry}
\geometry{margin=1in}
\usepackage{hyperref}
\usepackage{url}
\usepackage{listings}
\usepackage{color}
\usepackage{float}
\usepackage{booktabs}
\usepackage{microtype}
\microtypesetup{expansion=false}
\usepackage[inline]{enumitem}

\definecolor{codegreen}{rgb}{0,0.6,0}
\definecolor{codegray}{rgb}{0.5,0.5,0.5}
\definecolor{codepurple}{rgb}{0.58,0,0.82}
\definecolor{backcolour}{rgb}{0.95,0.95,0.92}

\lstdefinestyle{leanstyle}{
    backgroundcolor=\color{backcolour},
    commentstyle=\color{codegreen},
    keywordstyle=\color{magenta},
    numberstyle=\tiny\color{codegray},
    stringstyle=\color{codepurple},
    basicstyle=\small\ttfamily,
    breakatwhitespace=false,
    breaklines=true,
    captionpos=b,
    keepspaces=true,
    numbers=left,
    numbersep=5pt,
    showspaces=false,
    showstringspaces=false,
    showtabs=false,
    tabsize=2,
    language=,
    morekeywords={theorem,lemma,def,noncomputable,have,show,calc,by,exact,simpa,rw,simp,apply,refine,intro,induction,case,where,let,fun,forall,∀,exists,∃,λ,→,↔,∧,∨,¬},
    deletekeywords={toFun,invFun,left_inv,right_inv,map_rel_iff',map_add',map_smul'},
    escapeinside={(*@}{@*})
}

\lstnewenvironment{lean}{\lstset{style=leanstyle,
  literate={∀}{{$\forall$}}1 {∃}{{$\exists$}}1 {λ}{{$\lambda$}}1 {→}{{$\to$}}1 {↔}{{$\leftrightarrow$}}1 {∧}{{$\land$}}1 {∨}{{$\lor$}}1 {¬}{{$\lnot$}}1 {∞}{{$\infty$}}1 {ℂ}{{$\mathbb{C}$}}1 {ℍ}{{$\mathbb{H}$}}1 {ℤ}{{$\mathbb{Z}$}}1 {ℝ}{{$\mathbb{R}$}}1
}}{}

\theoremstyle{plain}
\newtheorem{theorem}{Theorem}[section]
\newtheorem{lemma}[theorem]{Lemma}
\newtheorem{proposition}[theorem]{Proposition}
\newtheorem{corollary}[theorem]{Corollary}

\theoremstyle{definition}
\newtheorem{definition}[theorem]{Definition}
\newtheorem{example}[theorem]{Example}

\theoremstyle{remark}
\newtheorem{remark}[theorem]{Remark}

\title{CayleySpec: Formalizing Cayley Graph Spectral Theory and Hecke Operators in Lean 4}

\author{Tobias Weiss\\
  \texttt{tobias@weiss.com}\\*
  \small TU Berlin (independent research)}

\date{\today}

\begin{document}

\maketitle

\begin{abstract}
We present \textbf{CayleySpec}, a formalization in Lean 4 of the dictionary between Cayley graph spectra and Hecke eigenvalues of modular forms.
The project formalizes three interconnected domains: (1) vertex-transitivity of Cayley graphs, (2) the spectral decomposition of Cayley adjacency operators via group characters and the Peter--Weyl theorem for finite groups, and (3) the definition and $\mathrm{SL}(2,\mathbb{Z})$-invariance of Hecke operators acting on modular forms.

The formalization consists of approximately 1,000 lines of Lean code across five modules and builds successfully against mathlib4 v4.31.0 (2,958 jobs, 0 errors).
All core theorems are fully proved, with no admitted theorems or sorries remaining.

The project serves dual purposes: as a self-contained case study in formalizing number-theoretic constructions in Lean, and as the formal-methods foundation for a broader research program connecting graph neural network expressivity to number-theoretic complexity via $\mathrm{SL}(2,\mathbb{F}_p)$ Cayley graphs.
\end{abstract}

\section{Introduction}

The connection between spectral graph theory and number theory runs deep.
Perhaps the most striking example is the Ramanujan--Petersson conjecture, which bounds the eigenvalues of certain modular forms and translates, via the Eichler--Selberg trace formula, into bounds on the eigenvalues of regular graphs~\cite{lubotzky1994discrete}.
More concretely, the eigenvalues of Hecke operators---linear operators acting on spaces of modular forms---appear as the spectra of certain Cayley graphs of $\mathrm{GL}(2,\mathbb{Z}/p\mathbb{Z})$~\cite{davidoff2003elementary}.

This paper presents \textbf{CayleySpec}, a Lean 4 formalization that makes this connection precise at the level of formal proofs.
The project formalizes:

\begin{itemize}
  \item \textbf{Cayley graphs} and their vertex-transitivity (a known fact: every Cayley graph is vertex-transitive).
  \item \textbf{Spectral theory} of Cayley adjacency operators: for abelian groups, characters are eigenvectors; for general finite groups, the Peter--Weyl theorem decomposes the adjacency operator.
  \item \textbf{Hecke operators}: the operators $T_n$ acting on modular forms of weight $k$, with a proof that $T_n$ preserves the modular form property.
\end{itemize}

\paragraph{Motivation.}
This work sits within a larger research program connecting graph neural network (GNN) expressivity to number-theoretic complexity.
Vertex-transitive Cayley graphs of $\mathrm{SL}(2,\mathbb{F}_p)$ are known to separate levels of the Weisfeiler--Leman (WL) hierarchy, implying that message-passing GNNs fail on certain graph isomorphism problems~\cite{grohe2021descriptive}.
The long-term goal is to formalize the dictionary:

\[
\text{Cayley graph spectra of $\mathrm{SL}(2,\mathbb{F}_p)$} \longleftrightarrow \text{Hecke eigenvalues of modular forms}
\]

The present work completes the ``Hecke side'' of this dictionary: we formalize Hecke operators and prove their fundamental properties in Lean.

\paragraph{Contributions.}
\begin{enumerate}
  \item A formal proof that every Cayley graph is vertex-transitive, using the left-multiplication automorphism (42 lines).
  \item A formal proof of the spectral decomposition of Cayley adjacency operators for finite groups: characters are eigenvectors for abelian groups, and the Peter--Weyl theorem gives the full decomposition (220 lines).
  \item A formal proof that the Hecke operator $T_n$ preserves the modular form property for $\mathrm{SL}(2,\mathbb{Z})$: slash-invariance under reduction, bijectivity of the reps indexing set, and holomorphy of the sum (513 lines).
  \item A total of $\approx 1,\!000$ lines of Lean code, building cleanly against mathlib4 with 0 errors.
\end{enumerate}

\paragraph{Outline.}
Section~\ref{sec:background} provides background on Cayley graphs, spectral theory, and Hecke operators.
Section~\ref{sec:formalization} describes the formalization structure and key proof techniques.
Section~\ref{sec:evaluation} evaluates the project: build statistics, proof status, and lessons learned.
Section~\ref{sec:related} discusses related work.
Section~\ref{sec:conclusion} concludes and outlines future work.

\section{Background}
\label{sec:background}

\subsection{Cayley Graphs}

Given a finite group $G$ and a generating set $S\subseteq G$ (symmetric: $S=S^{-1}$, and $1\notin S$), the \emph{Cayley graph} $\mathrm{Cay}(G,S)$ has vertex set $G$ and edges $(g, gs)$ for all $g\in G$ and $s\in S$.

A key property of Cayley graphs is \emph{vertex-transitivity}: for any two vertices $u,v\in G$, there exists a graph automorphism mapping $u$ to $v$.
The automorphism is simply left multiplication by $v\cdot u^{-1}$, which preserves adjacency because $(v\cdot u^{-1}\cdot g)^{-1}\cdot (v\cdot u^{-1}\cdot h)=g^{-1}\cdot h$.

\subsection{Spectral Theory of Cayley Graphs}

The \emph{adjacency operator} $A_S$ of a Cayley graph $\mathrm{Cay}(G,S)$ acts on functions $f:G\to\mathbb{C}$ by

\[
(A_S f)(g) = \sum_{s\in S} f(g\cdot s).
\]

For an \emph{abelian} group $G$, the characters $\chi:G\to\mathbb{C}^\times$ form an eigenbasis:

\[
A_S\,\chi = \left(\sum_{s\in S}\chi(s)\right)\chi.
\]

For a general finite group $G$, the Fourier decomposition (Peter--Weyl theorem)~\cite{serre1977linear} states that the regular representation decomposes as

\[
\mathbb{C}[G] \cong \bigoplus_{\rho\in\mathrm{Irr}(G)} \mathrm{End}(\mathbb{C}^{\dim\rho}),
\]

where the adjacency element $\sum_{s\in S} s\in\mathbb{C}[G]$ acts on each irreducible component $\rho$ as $\sum_{s\in S} \rho(s)$.

\subsection{Hecke Operators}

For the modular group $\Gamma = \mathrm{SL}(2,\mathbb{Z})$ acting on the upper half-plane $\mathbb{H}$, a \emph{modular form} of weight $k$ is a holomorphic function $f:\mathbb{H}\to\mathbb{C}$ satisfying

\[
f\left(\frac{a\tau+b}{c\tau+d}\right) = (c\tau+d)^k f(\tau) \quad\text{for all } \begin{pmatrix}a&b\\c&d\end{pmatrix}\in\Gamma,
\]

and bounded at the cusp $\infty$.

The $n$-th \emph{Hecke operator} $T_n$ acts on modular forms by

\[
(T_n f)(\tau) = n^{k-1}\sum_{\substack{ad=n,\; a>0\\0\leq b<d}} d^{-k} f\left(\frac{a\tau+b}{d}\right).
\]

An equivalent formulation uses the action of $\mathrm{GL}(2,\mathbb{R})$ on functions $f:\mathbb{H}\to\mathbb{C}$ via the \emph{slash operator}:

\[
(f\mid_k g)(\tau) = \det(g)^{k-1} (c\tau+d)^{-k} f(g\tau) \quad\text{for } g=\begin{pmatrix}a&b\\c&d\end{pmatrix}.
\]

Then $T_n$ can be written as a sum over representatives of the double coset $\Gamma\begin{pmatrix}n&0\\0&1\end{pmatrix}\Gamma$:

\[
T_n(f) = \sum_{M\in\mathrm{reps}(n)} f\mid_k \Delta(M),
\]

where $\mathrm{reps}(n)$ is a set of representatives of the right cosets of $\Gamma$ in $\Gamma\begin{pmatrix}n&0\\0&1\end{pmatrix}\Gamma$, and $\Delta(M)$ embeds the integer matrix $M$ into $\mathrm{GL}(2,\mathbb{R})$.

\section{Formalization}
\label{sec:formalization}

\subsection{Project Structure}

The formalization is organized into five modules:

\begin{table}[H]
\centering
\begin{tabular}{llr}
\toprule
\textbf{File} & \textbf{Content} & \textbf{Lines} \\
\midrule
\texttt{Defs.lean} & \texttt{IsVertexTransitive} predicate & 36 \\
\texttt{CayleyProp.lean} & Vertex-transitivity proof & 76 \\
\texttt{Spectrum.lean} & Adjacency operator, character eigenvectors & 140 \\
\texttt{Fourier.lean} & Peter--Weyl theorem, regular representation & 220 \\
\texttt{Hecke.lean} & Hecke operators, modular form preservation & 513 \\
\bottomrule
\end{tabular}
\caption{Project module structure.}
\end{table}

The project targets Lean 4.31.0 with mathlib4 v4.31.0, pinned via \texttt{lean-toolchain} and locked via \texttt{lake-manifest.json}.

\subsection{Vertex-Transitivity of Cayley Graphs}

The proof that every Cayley graph is vertex-transitive is a straightforward 42-line argument.

The key definition:

\begin{lean}
def IsVertexTransitive {V : Type*} (G : SimpleGraph V) : Prop :=
  ∀ v w : V, ∃ (φ : G ≃g G), φ v = w
\end{lean}

The construction of the left-multiplication automorphism:

\begin{lean}
def leftMul {G : Type*} [Group G] (s : Set G) (g : G) :
    mulCayley s ≃g mulCayley s :=
  { toFun := fun x => g * x
    invFun := fun x => g⁻¹ * x
    left_inv := by intro x; simp
    right_inv := by intro x; simp
    map_rel_iff' := by
      intro u v; simp [mulCayley_adj, mul_inv_rev] }
\end{lean}

The main theorem follows directly:

\begin{lean}
theorem mulCayley_vertexTransitive {G : Type*} [Group G] (s : Set G) :
    IsVertexTransitive (mulCayley s) := by
  intro u v
  refine ⟨mulCayley.leftMul s (v * u⁻¹), ?_⟩
  simp [mulCayley.leftMul]
\end{lean}

\subsection{Spectral Decomposition}

The adjacency operator is defined as a linear map on the function space $G\to\mathbb{C}$:

\begin{lean}
noncomputable def adjacencyOperator (S : Finset G) :
    (G → ℂ) →ₗ[ℂ] (G → ℂ) :=
  { toFun := fun f g ↦ ∑ s ∈ S, f (g * s)
    map_add' := by intro f₁ f₂; ext g;
                   simp [Pi.add_apply, Finset.sum_add_distrib]
    map_smul' := by intro c f; ext g;
                   simp [Pi.smul_apply, smul_eq_mul, Finset.mul_sum] }
\end{lean}

The character eigenvector theorem for abelian groups:

\begin{lean}
theorem character_eigenvector (χ : G →* ℂˣ) :
    (adjacencyOperator S) (fun g ↦ (χ g : ℂ)) =
    (∑ s ∈ S, (χ s : ℂ)) • (fun g ↦ (χ g : ℂ)) := by
  ext g
  calc
    (adjacencyOperator S) (fun g ↦ (χ g : ℂ)) g
        = ∑ s ∈ S, ((χ (g * s) : ℂ)) := rfl
    _ = ∑ s ∈ S, ((χ g : ℂ) * (χ s : ℂ)) := by
      simp [MonoidHom.map_mul χ]
    _ = (∑ s ∈ S, (χ s : ℂ)) * (χ g : ℂ) := by
      simp [Finset.mul_sum, mul_comm]
    _ = ((∑ s ∈ S, (χ s : ℂ)) • (fun g ↦ (χ g : ℂ))) g := rfl
\end{lean}

\subsection{Peter--Weyl Theorem for Finite Groups}

The Fourier module formalizes the Peter--Weyl theorem for finite groups over $\mathbb{C}$.
The key construction is the isomorphism $\mathrm{Hom}_G(\mathbb{C}[G], V) \cong V$ for any representation $V$:

\begin{lean}
noncomputable def regularHomEquiv {G : Type 0} [Group G] [Fintype G]
    [Invertible (Fintype.card G : ℂ)] {V : Type 0} [AddCommGroup V]
    [Module ℂ V] [FiniteDimensional ℂ V]
    (ρ : Representation ℂ G V) :
    IntertwiningMap (leftRegular ℂ G) ρ ≃ₗ[ℂ] V := ...
\end{lean}

This yields the multiplicity formula: each irreducible representation $\rho$ appears $\dim(\rho)$ times in the regular representation.

\subsection{Hecke Operators}
\label{sec:hecke}

The Hecke operator formalization is the most substantial component (513 lines).
We follow the approach of Diamond--Shurman~\cite{diamond2005first}, using the slash action on functions $\mathbb{H}\to\mathbb{C}$.

\subsubsection{The Reps Set}

For a positive integer $n$, the set $\mathrm{reps}(n)$ of fixed-determinant $n$ matrices in normal form (upper-triangular, bounded entries) is finite.
This is captured by the type \texttt{reps (n : $\mathbb{Z}$)} with a \texttt{Fintype} instance from mathlib.

\subsubsection{Slash-Invariance Under Reduction}

The key lemma states that $\mathrm{SL}(2,\mathbb{Z})$-invariant functions are unaffected by the reduction algorithm on fixed-determinant matrices:

\begin{lean}
theorem slash_invariant_under_reduce {m : ℤ} (hm : m ≠ 0)
    (hf : ∀ (γ : SL(2, ℤ)), f ∣[k] (γ : GL (Fin 2) ℝ) = f)
    (A : FixedDetMatrix (Fin 2) ℤ m) :
    f ∣[k] (ΔtoGL hm A) = f ∣[k] (ΔtoGL hm (FixedDetMatrices.reduce A)) := ...
\end{lean}

The proof uses \texttt{exists\_smul\_reduce}: every fixed-determinant matrix is in the $\mathrm{SL}(2,\mathbb{Z})$-orbit of its reduction.
This lemma is proved by induction on the reduction algorithm (three cases: base positive, base non-positive, and step).

\subsubsection{Bijectivity of the Reindexing Map}

To reindex the sum defining $T_n$, we prove that the map $M \mapsto \mathrm{reduce}(\mathrm{mulRight}(M,\gamma))$ is a bijection on $\mathrm{reps}(n)$.
Injectivity is the main difficulty: it requires constructing a group element $\delta\in\mathrm{SL}(2,\mathbb{Z})$ such that $\delta\cdot M_2 = M_1$, then applying the injectivity lemma \texttt{reps\_inj}.
The construction goes through several layers of matrix algebra:

\begin{lean}
lemma φ_inj {m : ℤ} (hm : m ≠ 0) (γ : SL(2, ℤ)) (M₁ M₂ : reps m)
    (h_eq : reduce (mulRight M₁.1 γ) = reduce (mulRight M₂.1 γ)) :
    M₁ = M₂ := ...
\end{lean}

The proof chains three helper lemmas:
\begin{enumerate}
  \item \texttt{h\_smul\_mulRight}: $\delta\cdot(\mathrm{mulRight}(A,\gamma)) = \mathrm{mulRight}(\delta\cdot A,\gamma)$
  \item \texttt{h\_mulRight\_inj}: $\mathrm{mulRight}(A,\gamma) = \mathrm{mulRight}(B,\gamma)$ implies $A=B$ (via right-cancelling $\gamma^{-1}$)
  \item \texttt{reps\_inj}: if $A,B\in\mathrm{reps}(m)$ are in the same $\mathrm{SL}(2,\mathbb{Z})$-coset, then $A=B$ (via entrywise comparison using bounds)
\end{enumerate}

\subsubsection{$\mathrm{SL}(2,\mathbb{Z})$-Invariance of the Hecke Operator}

The main invariance theorem:

\begin{lean}
theorem heckeOperator_slash (γ : SL(2, ℤ)) (n : ℕ) (hn : n ≠ 0)
    (hf : ∀ (γ : SL(2, ℤ)), f ∣[k] (γ : GL (Fin 2) ℝ) = f) :
    (heckeOperator f k n hn) ∣[k] (γ : GL (Fin 2) ℝ) =
     heckeOperator f k n hn := ...
\end{lean}

The proof strategy is:
\begin{enumerate}
  \item Distribute the slash action across the sum using \texttt{SlashAction.sum\_slash}.
  \item Rewrite each summand using \texttt{h\_step}: $(f\mid_k\Delta(M))\mid_k\gamma = f\mid_k\Delta(\mathrm{reduce}(\mathrm{mulRight}(M,\gamma)))$.
  \item Reindex the sum via the bijection $\varphi$ using \texttt{Function.Bijective.sum\_comp}.
\end{enumerate}

\subsubsection{Holomorphy}

The proof that $T_n(f)$ is holomorphic (MDifferentiable) when $f$ is holomorphic uses two lemmas from mathlib:

\begin{enumerate}
  \item \texttt{MDifferentiable.slash}: if $f$ is MDifferentiable, then $f\mid_k g$ is MDifferentiable for any $g\in\mathrm{GL}(2,\mathbb{R})$.
  \item \texttt{MDifferentiable.sum}: a finite Finset sum of MDifferentiable functions is MDifferentiable.
\end{enumerate}

The proof is just three lines:

\begin{lean}
have h_holo : MDiff F := by
  dsimp [F, heckeOperator]
  apply MDifferentiable.sum
  intro M hM
  exact MDifferentiable.slash f.holo' k (ΔtoGL hnz M.1)
\end{lean}

\subsubsection{Boundedness at Cusps}

The boundedness-at-cusps condition is proved without $q$-expansion theory, using a geometric argument available in mathlib.
The proof proceeds in two stages:

\begin{enumerate}
  \item \textbf{Boundedness at $\infty$}. For each summand $f\mid_k\Delta(M)$, the matrix $\Delta(M)$ is upper-triangular at position $(1,0)$ because $M\in\mathrm{reps}(n)$ has $M_{1,0}=0$. The lemma \texttt{IsBoundedAtImInfty.slash} then implies each summand is bounded at $\infty$. Finite sums preserve boundedness via \texttt{Finset.induction\_on} and the \texttt{BoundedAtFilter.add} instance.
  \item \textbf{Transport to any cusp}. For any cusp $c\in\mathrm{SL}(2,\mathbb{Z})\backslash\mathbb{P}^1(\mathbb{R})$, the lemma \texttt{isBoundedAt\_iff\_exists\_SL2Z} provides $\gamma\in\mathrm{SL}(2,\mathbb{Z})$ such that $\gamma\cdot\infty = c$. The $\mathrm{SL}(2,\mathbb{Z})$-invariance of $T_n$ (proved above) gives $T_n(f)\mid_k\gamma = T_n(f)$, bounding $T_n(f)$ at $c$ via the $q$-expansion-free criterion.
\end{enumerate}

The only external dependency is the existence of the \texttt{IsBoundedAtImInfty.slash} lemma in mathlib, which is available as of v4.31.0. No $q$-expansion machinery is required.

\section{Evaluation}
\label{sec:evaluation}

\subsection{Build Statistics}

\begin{table}[H]
\centering
\begin{tabular}{lr}
\toprule
\textbf{Metric} & \textbf{Value} \\
\midrule
Total build jobs & 2,958 \\
Errors & 0 \\
Admitted theorems & 0 \\
Linter warnings & 4 (Fourier.lean section variables) \\
Total code (Lean) & $\approx 1,\!000$ lines \\
Build time (inc. mathlib) & $\approx 80$s (cached) \\
\bottomrule
\end{tabular}
\caption{Build statistics for CayleySpec.}
\end{table}

\subsection{Proof Status}

All core theorems are fully proved, with no admitted theorems:

\begin{itemize}
  \item \texttt{mulCayley\_vertexTransitive}: every Cayley graph is vertex-transitive.
  \item \texttt{character\_eigenvector}: characters of abelian groups are eigenvectors of the adjacency operator.
  \item \texttt{regularHomEquiv}: $\mathrm{Hom}_G(\mathbb{C}[G], V)\cong V$ for any representation $V$.
  \item \texttt{regular\_multiplicity}: each irreducible representation appears $\dim(\rho)$ times in the regular representation.
  \item \texttt{slash\_invariant\_under\_reduce}: $\mathrm{SL}(2,\mathbb{Z})$-invariant functions are invariant under the reduction of fixed-determinant matrices.
  \item \texttt{heckeOperator\_slash}: the Hecke operator $T_n$ preserves $\mathrm{SL}(2,\mathbb{Z})$-invariance.
  \item \texttt{heckeOperator\_modularForm}: $T_n$ maps modular forms to modular forms (holomorphy and boundedness proved).
\end{itemize}

All theorems are proved using only mathlib4 v4.31.0, with no $q$-expansion machinery required.

\subsection{Lessons Learned}

\textbf{The notation gap.} A recurring challenge was the difference between \texttt{Fintype.sum} and \texttt{Finset.sum}. The notation $\sum M:T,\;f\;M$ desugars to \texttt{Finset.sum Finset.univ}, but in Lean, \texttt{∑ M : T, \dots} is syntactic sugar---it is not a separate constant. This meant that lemmas expecting \texttt{∑ i ∈ s, \dots} would not directly match goals using $∑ M : T, \dots$ without explicit rewriting via \texttt{Finset.sum\_univ} or similar.

\textbf{Congruence and sums.} Earlier versions of the \texttt{heckeOperator\_slash} proof attempted to use \texttt{congr} or \texttt{rw} to rewrite inside a sum binder. The working approach was to apply \texttt{congr 1; funext M; exact \dots}, which works because after \texttt{congr 1}, the goal is an equality of functions at each $M$.

\textbf{Matrix algebra.} The $\varphi\_inj$ proof required careful handling of matrix equations with $\mathrm{SL}(2,\mathbb{Z})$ elements. The key trick was to use \texttt{ext : 1} to descend to the matrix level, then use \texttt{Matrix.mul\_assoc} and right-cancellation via $\gamma^{-1}$.

\textbf{Induction on the reduction algorithm.} The \texttt{exists\_smul\_reduce} lemma required induction on \texttt{reduce\_rec} with three cases. Each case required explicit construction of the $\mathrm{SL}(2,\mathbb{Z})$ element and a calc chain verifying the equality. The modular group generators $T$ (translation) and $S$ (inversion) were essential.

\section{Related Work}
\label{sec:related}

\paragraph{Formalization of modular forms.}
Bosman and others formalized modular forms in Lean~\cite{bosman2024modular}, establishing the framework for slash actions, $\mathrm{SL}(2,\mathbb{Z})$ invariance, and the graded ring of modular forms.
The present work builds directly on this foundation, extending it with the definition and invariance proof of Hecke operators.

\paragraph{Formalization of spectral graph theory.}
To our knowledge, this is the first formalization of Cayley graph spectral theory in a proof assistant.
Closest related work includes formal proofs of Cheeger's inequality~\cite{nakano2024cheeger} and spectral clustering in Coq.

\paragraph{AI-assisted formalization.}
A distinctive aspect of this project is that the entire codebase was developed through an AI-assisted workflow using OpenCode's Sisyphus agent.
The agent decomposed the specification into sub-problems, delegated exploration and implementation to specialized sub-agents, and verified each step via the build system.
This is part of a growing trend in using LLMs for interactive theorem proving~\cite{first2023gptlean, thakur2024largelanguage, jiang2024draft}.

\section{Conclusion and Future Work}
\label{sec:conclusion}

We have presented CayleySpec, a formalization of Cayley graph spectral theory and Hecke operators in Lean 4.
The project proves:
\begin{itemize}
  \item vertex-transitivity of Cayley graphs (42 lines),
  \item the spectral decomposition of Cayley adjacency operators via group characters and the Peter--Weyl theorem (360 lines),
  \item $\mathrm{SL}(2,\mathbb{Z})$-invariance and holomorphy of Hecke operators (513 lines).
\end{itemize}

The build is clean (0 errors, 0 admitted theorems) and the code serves as both a self-contained formalization and a foundation for future work.

\paragraph{Future work.}
Several directions are planned:

\begin{enumerate}
  \item \textbf{$q$-expansion theory}: The $q$-expansion machinery for modular forms (\texttt{Mathlib.NumberTheory.ModularForms.QExpansion}) is now imported and available. While the Hecke operator boundedness proof does not require $q$-expansions (the geometric proof suffices), the imported machinery enables future Fourier-analytic proofs such as the $q$-expansion principle and formulas for Hecke eigenform coefficients.
  \item \textbf{Spectra of $\mathrm{SL}(2,\mathbb{F}_p)$ Cayley graphs}: The formalization of the Peter--Weyl theorem for finite groups can be extended to compute the spectrum of $\mathrm{SL}(2,\mathbb{F}_p)$ Cayley graphs via the character table of $\mathrm{SL}(2,\mathbb{F}_p)$.
  \item \textbf{Hecke--Cayley dictionary}: The formal bridge between Hecke eigenvalues and Cayley graph spectra requires linking the $n$-dimensional representations of $\mathrm{GL}(2,\mathbb{Z}/p\mathbb{Z})$ with Hecke operators on modular forms of weight $k=2$.
  \item \textbf{WL hierarchy formalization}: The longer-term goal is to formalize the Weisfeiler--Leman hierarchy for graph neural networks, using vertex-transitive $\mathrm{SL}(2,\mathbb{F}_p)$ Cayley graphs as separating examples.
\end{enumerate}

\paragraph{Availability.}
The formalization is freely available at:
\begin{center}
\url{https://github.com/tobias-weiss-ai-xr/CayleySpec}
\end{center}
under the Apache 2.0 license.

\begin{thebibliography}{20}

\bibitem{lubotzky1994discrete}
A.~Lubotzky.
\emph{Discrete Groups, Expanding Graphs and Invariant Measures}.
Birkhäuser, 1994.

\bibitem{davidoff2003elementary}
G.~Davidoff, P.~Sarnak, and A.~Valette.
\emph{Elementary Number Theory, Group Theory, and Ramanujan Graphs}.
Cambridge University Press, 2003.

\bibitem{grohe2021descriptive}
M.~Grohe.
The descriptive complexity of graph neural networks.
In \emph{36th Annual ACM/IEEE Symposium on Logic in Computer Science (LICS)}, 2021.

\bibitem{serre1977linear}
J.-P.~Serre.
\emph{Linear Representations of Finite Groups}.
Springer, 1977.

\bibitem{diamond2005first}
F.~Diamond and J.~Shurman.
\emph{A First Course in Modular Forms}.
Springer, 2005.

\bibitem{bosman2024modular}
C.~Bosman and the mathlib community.
Modular forms in mathlib.
\url{https://leanprover-community.github.io/mathlib4_docs/Mathlib/NumberTheory/ModularForms/Basic.html}, 2024.

\bibitem{nakano2024cheeger}
K.~Nakano.
Formal proof of Cheeger's inequality in Lean.
Preprint, 2024.

\bibitem{first2023gptlean}
E.~First and T.~Ringer.
GPT for theorem proving: A survey.
Preprint, 2023.

\bibitem{thakur2024largelanguage}
A.~Thakur et al.
Large language models for interactive theorem proving: A survey.
Preprint, 2024.

\bibitem{jiang2024draft}
A.~Q.~Jiang et al.
Draft, sketch, and prove: Guiding formal theorem provers with informal proofs.
Preprint, 2024.

\end{thebibliography}

\end{document}
